\documentclass[12pt,a4paper]{article}
\usepackage[utf8]{inputenc}
\usepackage[margin=1in]{geometry}
\usepackage{enumitem}
\usepackage{fancyhdr}
\usepackage{titlesec}
\usepackage{xcolor}
\usepackage{tcolorbox}
\usepackage[hidelinks, unicode=false]{hyperref}
\usepackage{parskip}
\usepackage{mdframed}
\usepackage{array}

\setlength{\headheight}{25pt}
\addtolength{\topmargin}{-13pt}

\definecolor{awsorange}{RGB}{255,153,0}
\definecolor{awsblue}{RGB}{35,47,62}
\definecolor{lightgray}{RGB}{245,245,245}
\definecolor{analogyblue}{RGB}{232,244,255}

\pagestyle{fancy}
\fancyhf{}
\fancyhead[L]{\textcolor{awsblue}{\textbf{Cloud Operations 2 -- Course Overview}}}
\fancyhead[R]{\textcolor{awsorange}{\textbf{ACI Q2}}}
\fancyfoot[C]{\thepage}
\renewcommand{\headrulewidth}{2pt}
\renewcommand{\headrule}{%
  \hbox to\headwidth{\color{awsorange}\leaders\hrule height\headrulewidth\hfill}}

\titleformat{\section}{\large\bfseries\color{awsblue}}{}{0em}{}[%
  \color{awsorange}\titlerule]
\titlespacing{\section}{0pt}{16pt}{6pt}

\newmdenv[
  backgroundcolor=analogyblue,
  linecolor=awsorange,
  linewidth=1.5pt,
  leftline=true, rightline=false, topline=false, bottomline=false,
  innerleftmargin=10pt, innerrightmargin=10pt,
  innertopmargin=6pt,  innerbottommargin=6pt
]{analogybox}

\begin{document}

% --------------------------------------------------------
%  TITLE PAGE
% --------------------------------------------------------
\begin{titlepage}
  \centering
  \vspace*{3cm}
  {\fontsize{34}{42}\selectfont\bfseries\color{awsblue}Cloud Operations 2\par}
  \vspace{0.5cm}
  {\fontsize{18}{26}\selectfont\color{awsorange}Course Overview \& Weekly Guide\par}
  \vspace{0.5cm}
  {\large AWS Cloud Institute -- Q2\par}
  \vspace{2.5cm}
  \begin{tcolorbox}[
    colback=lightgray, colframe=awsblue, width=0.85\textwidth, arc=6pt,
    title={\textbf{About This Document}},
    fonttitle=\color{white}, colbacktitle=awsblue]
  This document provides a concise overview of all 11 weeks of Cloud Operations 2.
  Each week includes the core topics, key AWS services, and a short analogy to anchor
  the concept before you dive into the full module.

  \vspace{0.4cm}
  \textbf{Course themes:} DevOps \textbullet{} Infrastructure as Code \textbullet{}
  Monitoring \textbullet{} Troubleshooting \textbullet{} Distributed Tracing
  \end{tcolorbox}
  \vspace{2cm}
  {\large \today\par}
\end{titlepage}

\tableofcontents
\newpage

% --------------------------------------------------------
%  COURSE OVERVIEW
% --------------------------------------------------------
\section*{Course Overview}
\addcontentsline{toc}{section}{Course Overview}

Cloud Operations 2 builds on foundational AWS knowledge and focuses on the
\textbf{operational side} of running software in the cloud.
Over 11 weeks you will move through three broad themes:

\begin{enumerate}
  \item \textbf{Weeks 1--3: DevOps} -- automating how code is built, tested,
        deployed, and secured.
  \item \textbf{Weeks 4--5: Infrastructure as Code} -- treating your infrastructure
        the same way you treat application code.
  \item \textbf{Weeks 6--11: Observability} -- scaling, logging, monitoring, and
        tracing so you always know what your system is doing.
\end{enumerate}

\textbf{Completion requirement:} Score 80\% or higher on the final assessment at the
end of each module. Assessments can be retaken unlimited times.

\newpage

% --------------------------------------------------------
%  WEEK 1
% --------------------------------------------------------
\section{Week 1 -- DevOps Part 1: CI/CD and Automated Testing}

\textbf{Goal:} Understand the principles of Continuous Integration and Continuous
Delivery (CI/CD), automate testing with AWS CodeBuild, and integrate code reviews
into the pipeline.

\textbf{Key topics:}
\begin{itemize}[leftmargin=1.5em]
  \item CI/CD core principles, benefits, and best practices
  \item Automating unit testing -- code coverage and why it matters
  \item AWS CodeBuild: types of tests it can run and how to configure them
  \item Integrating code reviews with Amazon CodeGuru
\end{itemize}

\textbf{Key services:} AWS CodePipeline, AWS CodeBuild, Amazon CodeGuru

\begin{analogybox}
\textbf{Analogy -- The Assembly Line Quality Check}

Think of a car assembly line. Workers do not wait until the car is fully built to
check for defects -- they inspect each part as it is added. CI/CD works the same
way: every time a developer pushes code, automated tests run immediately, catching
defects early before they reach the customer.
\end{analogybox}

\textbf{Labs:} CI/CD Pipelines for APIs (SimuLearn) \textbullet{}
Unit Testing with CodePipeline \textbullet{}
Automating Code Reviews with CodeGuru

\newpage

% --------------------------------------------------------
%  WEEK 2
% --------------------------------------------------------
\section{Week 2 -- DevOps Part 2: Deployment Strategies and DevOps Environments}

\textbf{Goal:} Learn how CodePipeline and CodeDeploy work together to deliver
software, explore deployment strategies, and understand how to standardize DevOps
environments with AWS Proton and AWS Service Catalog.

\textbf{Key topics:}
\begin{itemize}[leftmargin=1.5em]
  \item The developer's role in continuous delivery
  \item Deployment strategies: In-Place, Rolling, Immutable, Blue/Green
  \item AWS CodeDeploy components and configuration
  \item AWS Proton for environment standardization
  \item AWS Service Catalog for self-service infrastructure
\end{itemize}

\textbf{Key services:} AWS CodePipeline, AWS CodeDeploy, AWS Proton, AWS Service Catalog

\begin{analogybox}
\textbf{Analogy -- Painting a Bridge}

A Blue/Green deployment is like painting a bridge. You do not close the bridge and
repaint it while traffic is stopped. Instead, you build a second bridge next to it,
paint that one, then redirect traffic. If anything is wrong, you switch back
instantly -- the old bridge is still there.
\end{analogybox}

\textbf{Labs:} Integration Testing with CodePipeline \textbullet{}
Blue/Green Deployments with CodeDeploy and EC2 \textbullet{}
Deploying Standardized Assets with Service Catalog

\newpage

% --------------------------------------------------------
%  WEEK 3
% --------------------------------------------------------
\section{Week 3 -- DevOps Part 3: DevSecOps}

\textbf{Goal:} Understand DevSecOps and how to integrate security practices
throughout the software development lifecycle rather than treating security as an
afterthought.

\textbf{Key topics:}
\begin{itemize}[leftmargin=1.5em]
  \item The role and key components of DevSecOps
  \item Common software vulnerabilities (OWASP, CVEs, misconfigurations)
  \item Security testing practices: SAST, DAST, dependency scanning
  \item Amazon CodeGuru Security and Amazon Inspector
  \item Automating vulnerability management at scale
\end{itemize}

\textbf{Key services:} Amazon Inspector, Amazon CodeGuru Security, AWS Security Hub

\begin{analogybox}
\textbf{Analogy -- Smoke Detectors vs.\ Fire Trucks}

Traditional security is like calling the fire truck after the building is already
burning. DevSecOps is like installing smoke detectors in every room -- you catch
the problem at the earliest possible moment, when it is cheapest and easiest to fix.
\end{analogybox}

\textbf{Labs:} Vulnerability Scanning with Amazon Inspector \textbullet{}
Secure Self-Service Infrastructure (SimuLearn) \textbullet{}
Securing Your Servers (SimuLearn)

\newpage

% --------------------------------------------------------
%  WEEK 4
% --------------------------------------------------------
\section{Week 4 -- Infrastructure as Code Part 1: Foundations and CloudFormation}

\textbf{Goal:} Understand what Infrastructure as Code (IaC) is, how it fits into a
DevOps workflow, and how to build reusable, parameterized AWS CloudFormation templates.

\textbf{Key topics:}
\begin{itemize}[leftmargin=1.5em]
  \item What IaC is and the problems it solves over manual provisioning
  \item Developer and infrastructure team responsibilities in IaC
  \item AWS and non-AWS IaC tools (CloudFormation, CDK, Terraform)
  \item CloudFormation template anatomy: intrinsic functions, conditions,
        rules, DependsOn, macros, and modules
  \item CloudFormation Designer
\end{itemize}

\textbf{Key services:} AWS CloudFormation, AWS CDK, AWS CodePipeline

\begin{analogybox}
\textbf{Analogy -- Architectural Blueprints}

Before a building is constructed, architects produce detailed blueprints. Anyone
with those blueprints can build an identical structure anywhere in the world.
CloudFormation templates are your cloud blueprints -- version-controlled,
repeatable, and shareable across every environment.
\end{analogybox}

\textbf{Labs:} Deploying IaC with CodePipeline \textbullet{}
Building a Serverless API with Infrastructure Composer \textbullet{}
CloudFormation Intrinsic Functions and Keywords

\newpage

% --------------------------------------------------------
%  WEEK 5
% --------------------------------------------------------
\section{Week 5 -- Infrastructure as Code Part 2: Advanced Stacks and AWS CDK}

\textbf{Goal:} Go deeper into CloudFormation stacks (nested stacks, drift detection,
failure handling, change sets) and introduce the AWS Cloud Development Kit (CDK) as a
code-first alternative to YAML templates.

\textbf{Key topics:}
\begin{itemize}[leftmargin=1.5em]
  \item Stacks and change sets -- building, updating, and rolling back
  \item Stack drift detection and remediation
  \item Nested stacks and deploying multiple stacks
  \item CloudFormation Registry and EventBridge integration
  \item AWS CDK: provisioning resources with familiar programming languages
  \item Testing CDK constructs
\end{itemize}

\textbf{Key services:} AWS CloudFormation, AWS CDK, Amazon EventBridge

\begin{analogybox}
\textbf{Analogy -- LEGO vs.\ Custom Bricks}

CloudFormation YAML is like standard LEGO bricks -- reliable and well-documented.
AWS CDK is like writing a program that \emph{generates} the LEGO instructions for
you, using loops, conditions, and reusable functions. Both produce the same result,
but CDK lets developers work in the language they already know.
\end{analogybox}

\textbf{Labs:} Deploying Advanced CloudFormation Stacks \textbullet{}
Automation with CloudFormation (SimuLearn) \textbullet{}
Working with Constructs in AWS CDK

\newpage

% --------------------------------------------------------
%  WEEK 6
% --------------------------------------------------------
\section{Week 6 -- Logging and Scaling Part 1: Auto Scaling and Load Balancing}

\textbf{Goal:} Learn how AWS Auto Scaling keeps your application right-sized under
varying load, and how Elastic Load Balancing distributes traffic for high availability
and zero-downtime deployments.

\textbf{Key topics:}
\begin{itemize}[leftmargin=1.5em]
  \item How Auto Scaling works across EC2, ECS, DynamoDB, and Aurora
  \item Scaling policies: target tracking, step scaling, scheduled, predictive
  \item Launch templates, Spot Fleets, lifecycle hooks
  \item Elastic Load Balancing types (ALB, NLB) and how they work
  \item Using ALB for Blue/Green and canary deployments
  \item ELB integration with Route~53 and AWS PrivateLink
\end{itemize}

\textbf{Key services:} Amazon EC2 Auto Scaling, Elastic Load Balancing, Amazon Route~53

\begin{analogybox}
\textbf{Analogy -- A Restaurant on a Busy Friday Night}

Auto Scaling is like a restaurant that calls in extra staff when the dinner rush
hits and sends them home when it quiets down. The load balancer is the host at the
door -- directing each customer to a table so no single waiter is overwhelmed while
others stand idle.
\end{analogybox}

\textbf{Labs:} Introduction to EC2 Auto Scaling \textbullet{}
Introduction to Elastic Load Balancing \textbullet{}
Highly Available Web Applications (SimuLearn)

\newpage

% --------------------------------------------------------
%  WEEK 7
% --------------------------------------------------------
\section{Week 7 -- Logging and Scaling Part 2: CloudWatch Monitoring}

\textbf{Goal:} Build a comprehensive monitoring strategy using Amazon CloudWatch --
from basic metrics and dashboards to anomaly detection, container insights, and
security monitoring.

\textbf{Key topics:}
\begin{itemize}[leftmargin=1.5em]
  \item Logging overview and monitoring strategy
  \item CloudWatch metrics, dashboards, and the CloudWatch Agent
  \item AWS CloudTrail for API activity logging
  \item Serverless observability, Container Insights, Application Insights
  \item CloudWatch Anomaly Detection and Metric Streams
  \item CloudWatch Synthetics, RUM, Logs Insights, Contributor Insights
  \item High-resolution metrics, composite alarms, and scaling with alarms
  \item Introduction to AWS X-Ray
\end{itemize}

\textbf{Key services:} Amazon CloudWatch, AWS CloudTrail, AWS X-Ray

\begin{analogybox}
\textbf{Analogy -- The Dashboard of a Car}

You do not wait for your car to break down to check its health. The dashboard shows
speed, fuel, temperature, and warning lights in real time. CloudWatch is your
application's dashboard -- it tells you what is happening right now and warns you
before something breaks.
\end{analogybox}

\textbf{Labs:} Collecting and Analyzing Logs with CloudWatch Logs Insights
\textbullet{} Monitoring Security with Amazon CloudWatch

\newpage

% --------------------------------------------------------
%  WEEK 8
% --------------------------------------------------------
\section{Week 8 -- Monitoring and Troubleshooting Part 1: CloudTrail and IAM}

\textbf{Goal:} Apply a structured troubleshooting methodology to cloud infrastructure
problems, use CloudTrail to investigate security incidents, and diagnose common IAM
permission issues.

\textbf{Key topics:}
\begin{itemize}[leftmargin=1.5em]
  \item Structured troubleshooting methodology (define, gather, hypothesize, test, fix)
  \item CloudTrail event types and log file analysis
  \item CloudTrail integration with other AWS services and CloudTrail Lake
  \item Troubleshooting IAM policy errors and role assumption issues
  \item Trust relationships and service-linked roles
  \item Best practices for IAM policy management
\end{itemize}

\textbf{Key services:} AWS CloudTrail, AWS IAM, CloudTrail Lake

\begin{analogybox}
\textbf{Analogy -- The Black Box Recorder}

Every commercial aircraft carries a black box that records every action taken in
the cockpit. CloudTrail is your AWS black box -- every API call, who made it,
when, and from where. When something goes wrong, it is the first place investigators
look.
\end{analogybox}

\textbf{Labs:} Monitoring and Troubleshooting EC2 Workloads with Detective Controls
\textbullet{} Performing a Basic Audit of Your AWS Environment

\newpage

% --------------------------------------------------------
%  WEEK 9
% --------------------------------------------------------
\section{Week 9 -- Monitoring and Troubleshooting Part 2: Advanced Metrics}

\textbf{Goal:} Go beyond standard metrics with custom metrics, cross-account
monitoring, anomaly detection, and purpose-built insights tools for containers,
Lambda, and networks.

\textbf{Key topics:}
\begin{itemize}[leftmargin=1.5em]
  \item Publishing and using custom CloudWatch metrics
  \item Cross-account, cross-Region monitoring
  \item CloudWatch Anomaly Detection and Metrics Insights alarms
  \item CloudWatch Container Insights for ECS and EKS
  \item CloudWatch Lambda Insights -- cold starts, memory, duration
  \item CloudWatch Internet Monitor and Network Monitor
\end{itemize}

\textbf{Key services:} Amazon CloudWatch -- Container Insights, Lambda Insights,
Internet Monitor, Network Monitor

\begin{analogybox}
\textbf{Analogy -- Specialist Doctors vs.\ a General Practitioner}

Standard CloudWatch metrics are like a GP -- great for a general health check.
Container Insights and Lambda Insights are specialists: they know exactly which
questions to ask about containers and serverless functions, surfacing details a
general check would miss.
\end{analogybox}

\textbf{Labs:} Monitor and Analyze Network Traffic (SimuLearn)
\textbullet{} Traffic Mirroring (SimuLearn)

\newpage

% --------------------------------------------------------
%  WEEK 10
% --------------------------------------------------------
\section{Week 10 -- Monitoring and Troubleshooting Part 3: Log Analysis}

\textbf{Goal:} Master log-based monitoring using CloudWatch Logs -- metric filters,
Live Tail, Logs Insights queries, anomaly detection, VPC Flow Logs troubleshooting,
and application-level monitoring tools.

\textbf{Key topics:}
\begin{itemize}[leftmargin=1.5em]
  \item Metric filters -- extracting metrics from log data in real time
  \item Live Tail -- streaming logs during deployments and incidents
  \item CloudWatch Logs Insights -- querying and analyzing log data
  \item Log anomaly detection and pattern analysis
  \item Troubleshooting with VPC Flow Logs
  \item CloudWatch Application Insights, Synthetics, RUM, Application Signals
\end{itemize}

\textbf{Key services:} Amazon CloudWatch Logs, VPC Flow Logs, CloudWatch Synthetics,
CloudWatch RUM

\begin{analogybox}
\textbf{Analogy -- CCTV Footage}

Logs are like CCTV footage of your application. Metric filters are motion-detection
alerts that trigger when something specific happens. Live Tail is watching the live
feed during an incident. Logs Insights is rewinding and searching the recording to
find exactly what happened and when.
\end{analogybox}

\textbf{Labs:} Monitoring Applications and Infrastructure

\newpage

% --------------------------------------------------------
%  WEEK 11
% --------------------------------------------------------
\section{Week 11 -- Monitoring and Troubleshooting Part 4: Distributed Tracing}

\textbf{Goal:} Understand distributed tracing and use AWS X-Ray to follow a single
request as it travels through multiple services, identify latency bottlenecks, and
troubleshoot network reachability issues.

\textbf{Key topics:}
\begin{itemize}[leftmargin=1.5em]
  \item X-Ray concepts: traces, segments, subsegments, annotations, metadata
  \item Service graphs and trace maps
  \item Filter expressions and latency histograms
  \item X-Ray Insights and the Analytics console
  \item Cross-account tracing and X-Ray SDK for Python
  \item VPC Reachability Analyzer for network troubleshooting
\end{itemize}

\textbf{Key services:} AWS X-Ray, VPC Reachability Analyzer, CloudWatch ServiceLens

\begin{analogybox}
\textbf{Analogy -- Tracking a Package}

When you order something online, the tracking system shows every stop the package
made: warehouse, sorting facility, delivery truck, your door. X-Ray does the same
for a web request -- it shows every service the request touched, how long each step
took, and exactly where it slowed down or failed.
\end{analogybox}

\textbf{Labs:} Troubleshooting Website Reachability Behind a Load Balancer
\textbullet{} Core Security Concepts (SimuLearn)

\newpage

% --------------------------------------------------------
%  QUICK REFERENCE TABLE
% --------------------------------------------------------
\section*{Quick Reference -- Services by Week}
\addcontentsline{toc}{section}{Quick Reference -- Services by Week}

\begin{center}
\renewcommand{\arraystretch}{1.5}
\begin{tabular}{|p{0.07\textwidth}|p{0.37\textwidth}|p{0.43\textwidth}|}
\hline
\textbf{Week} & \textbf{Theme} & \textbf{Key AWS Services} \\
\hline
1  & CI/CD and Automated Testing            & CodePipeline, CodeBuild, CodeGuru \\
\hline
2  & Deployment Strategies and Environments & CodeDeploy, Proton, Service Catalog \\
\hline
3  & DevSecOps                              & Inspector, CodeGuru Security, Security Hub \\
\hline
4  & IaC Foundations and CloudFormation     & CloudFormation, CDK, CodePipeline \\
\hline
5  & Advanced Stacks and CDK                & CloudFormation, CDK, EventBridge \\
\hline
6  & Auto Scaling and Load Balancing        & EC2 Auto Scaling, ELB, Route~53 \\
\hline
7  & CloudWatch Monitoring                  & CloudWatch, CloudTrail, X-Ray \\
\hline
8  & Troubleshooting: CloudTrail and IAM    & CloudTrail, IAM, CloudTrail Lake \\
\hline
9  & Advanced Metrics and Insights          & CloudWatch Insights, Internet Monitor \\
\hline
10 & Log Analysis and App Monitoring        & CloudWatch Logs, VPC Flow Logs, RUM \\
\hline
11 & Distributed Tracing with X-Ray         & X-Ray, Reachability Analyzer, ServiceLens \\
\hline
\end{tabular}
\end{center}

\vspace{1.5cm}

\begin{tcolorbox}[
  colback=lightgray, colframe=awsorange, arc=6pt,
  title={\textbf{Exam Alignment -- SysOps Administrator Associate (SOA-C02)}},
  fonttitle=\color{white}, colbacktitle=awsorange]
\begin{itemize}[leftmargin=1.5em, itemsep=2pt]
  \item \textbf{Domain 1} Monitoring, Logging, Remediation (20\%) $\rightarrow$ Weeks 7--11
  \item \textbf{Domain 2} Reliability and Business Continuity (16\%) $\rightarrow$ Week 6
  \item \textbf{Domain 3} Deployment, Provisioning, Automation (18\%) $\rightarrow$ Weeks 1--2, 4--5
  \item \textbf{Domain 4} Security and Compliance (16\%) $\rightarrow$ Weeks 3, 8
  \item \textbf{Domain 5} Networking and Content Delivery (18\%) $\rightarrow$ Weeks 6, 11
  \item \textbf{Domain 6} Cost and Performance Optimization (12\%) $\rightarrow$ Weeks 6, 9
\end{itemize}
\end{tcolorbox}

\end{document}
